\clearpage
\fancyhead[L]{\small BODE. ASTRONOMISCHES JAHRBUCH.}
\SourceLeaf{509}
\begin{center}\small Göttingische gelehrte Anzeigen. Stück 155. Seite 1545..1552. 1808. September 26.\\[3pt]\rule{.72\linewidth}{.3pt}\end{center}

Bei dem Verfasser, und in Commission bei \textsc{Fr. Braunes}: \emph{Astronomisches Jahrbuch für das Jahr} 1810, \emph{nebst einer Sammlung der neuesten in die astronomischen Wissenschaften einschlagenden Abhandlungen und Nachrichten. Von} \textsc{J. E. Bode}. 1807. 268 Seiten in Octav, und eine Kupfertafel. Berlin.

Die Einrichtung des ersten Theils dieses geschätzten und jedem Astronomen unentbehrlichen Jahrbuches, der den astronomischen Kalender enthält, ist dieselbe, wie in dem vorhergehenden Jahrgange. Da indess der würdige Hr. Herausgeber von jeher unermüdet die Brauchbarkeit desselben zu erhöhen bemühet gewesen ist, und dahin abzweckende Vorschläge und Wünsche nicht unbeachtet gelassen hat; so erlauben wir uns hier einige solche Wünsche darzulegen, die, wie wir bestimmt wissen, auch mehrere andere Astronomen längst gehegt haben. Die ausführliche Erläuterung des Kalenders und seines Gebrauchs ist seit dem Jahrgange 1800 nicht wieder abgedruckt; wir finden dies sehr lobenswerth, da dieser ziemlich starke Aufsatz einen besser zu nützenden Platz einnehmen würde, und man ausserdem \emph{das Meiste}, was er enthält, ohnedies bei jedem als bekannt voraussetzen darf, der sich mit Beobachtung des Himmels beschäftigen will. Allein Einiges aus dieser Erklärung, was nemlich das Willkührliche und sich nicht von selbst Verstehende bei dem Kalender bestrifft, sollte doch wirklich zum stehenden Artikel gemacht, oder sonst auf irgend eine Art dem Kalender einverleibt werden. Dass z. B. die Planetenörter alle für Mitternacht wahre Berliner Zeit angesetzt sind, mit Ausnahme der Venus und des Merkur, deren Stellungen für den Mittag gelten, ist Etwas, das Niemand von selbst erräth, und worüber Jeder in Ungewissheit bleibt, der nicht Gelegenheit hat, den Jahrgang von 1800 nachzuschlagen. Eben so, dass die Bedeckungen der Fixsterne vom Monde für wahre, und nicht für die jetzt bei fast allen Astronomen gebräuchliche mittlere Zeit gelten. Bei der scheinbaren Schiefe der Ekliptik wird die Nutation immer mit verändertem Zeichen angegeben: ohne Erklärung versteht sich dies doch nicht von selbst, und man sollte eher glauben, dass man, da doch die scheinbare Schiefe, die mit der Nutation behaftet ist, von jener diese abziehen müsste, um die mittlere zu erhalten, da es sich doch mit der im Jahrbuch angesetzten umgekehrt verhält. Auch können wir es nicht billigen, dass Hr. \textsc{Bode} die Schiefe der Ekliptik noch immer nach Hrn. \textsc{von Zach}'s ältern Sonnentafeln angibt, die nach den neuern, schon seit vielen Jahren von Englischen, Französischen und Italiänischen Astronomen angestellten, Beobachtungen um $8''$ zu gross ist, daher auch alle Declinationen der Sonne im Jahrbuche zu gross sind, und also nach Verschiedenheit der Jahreszeiten unrichtige Polhöhen geben, wenn man sich auf dieselben verlässt. Zu wünschen wäre es auch, dass Hr. \textsc{Bode} die unbequeme Angabe der Zeit nach Vormittags- und Nachmittagsstunden mit der bei den Astronomen üblichen vertauschen möchte. Es sind uns mehrere Beispiele bekannt, dass Sternbedeckungen und Jupiterstrabanten-Verfinsterungen blos wegen der Angabe nach bürgerlicher Zeit verfehlt sind. Bei den Stellungen der Planeten hat Hr. \textsc{Bode} nach einem Wunsche, den ein Recensent im 41. Stück der Literaturzeitung 1797 geäussert hatte, seit dem J. 1800 auch die heliocentrischen Längen und Breiten aufgenommen, aber dafür die geraden Aufsteigungen weggelassen: wir sind damit nie zufrieden gewesen, und glauben, dass die geocentrische gerade Aufsteigung und Abweichung das Allerwesentlichste in einer astronomischen Ephemeride ist; wir würden lieber alle übrigen sich auf die Planeten beziehenden Columnen entbehren, als diese beiden nothwendigen Stücke; lieb würde es gewiss in mancher Rücksicht jedem Astronomen sein, wenn die Abstände von der Erde beigefügt würden, die zur Bestimmung der Aberration unentbehrlich sind. Die Gründe, warum die heliocentrischen Oerter von jenem Rec. gewünscht wurden, scheinen uns so erheb\SourceLeaf{510}lich nicht, und wir glauben, dass von diesen beiden Columnen wenig Gebrauch gemacht werden kann, zumal da Hr. \textsc{Bode} unter den monatlichen Erscheinungen die Oppositionen und Conjunctionen der Planeten angibt, welche sich allenfalls eben so leicht aus den geocentrischen ableiten lassen. Wir könnten hier noch manche andre Wünsche beifügen, die ein vieljähriger Gebrauch der astronomischen Ephemeriden veranlasst hat; da sich indessen dieselben nicht wohl ohne grössere Abänderungen in der Gestalt und Einrichtung derselben ausführen liessen, so schränken wir uns hier nur auf einen ein, durch dessen Erfüllung der verdienstvolle Herausgeber alle diejenigen sehr verpflichten würde, die viel mit Planeten- und Cometenrechnungen zu thun haben, und denen die Berechnung von Sonnenörtern immer einen sehr grossen Zeitaufwand verursacht. Dieser könnte in sehr vielen Fällen ganz vermieden werden, wenn die Abstände der Sonne von der Erde, oder deren Logarithmen, allenfalls sogar mit Einer Decimalstelle weniger, als Hr. \textsc{Bode} in den vier letzten Jahrgängen geliefert hat, für alle einzelne Tage des Jahrs gegeben würden; für die dadurch etwas vergrösserte Arbeit könnte der Herausgeber sich durch Weglassung der Secunden bei den Längen und Breiten des Mondes schadlos halten, deren Berechnung eine ungeheure Arbeit kostet, und von denen, wie wir glauben behaupten zu können, doch wenig oder gar kein Gebrauch gemacht wird.

Nach diesen Bemerkungen über die Einrichtung des Jahrbuchs, zu denen uns blos unser aufrichtiger Wunsch, die Brauchbarkeit desselben noch erhöhet zu sehen, veranlasst hat, wenden wir uns zu den Abhandlungen, wodurch das astronomische Jahrbuch einen dauernden Werth erhält. Diese sind: 1) Untersuchung der wahren elliptischen Bewegung des Cometen von 1769, von \textsc{Fr. W. Bessel}. Diese Schrift wurde zur Concurrenz um einen Preis von 30 Friedrichsd'or eingesandt, welchen ein Ungenannter auf die beste astronomische Abhandlung oder die merkwürdigste Entdeckung gesetzt hatte; man erkannte aber Hrn. \textsc{Bessel} nur die eine Hälfte zu, und die andere Hrn. \textsc{Huth}, als Mitentdecker der beiden Cometen des J. 1805. Die elliptische Bahn jenes Cometen hatte schon mehrere Astronomen beschäftigt, unter denen der Pater \textsc{Asklepi} diese Untersuchung am sorgfältigsten behandelt hatte. Indess verdiente dieser Gegenstand allerdings eine neue Bearbeitung, wozu besonders die grossen Reformen unsrer Sternverzeichnisse und Sonnentafeln einluden, mit deren Hülfe man viel zuverlässigere Resultate zu erhalten hoffen durfte. Hr. \textsc{Bessel} hat ihn jetzt so musterhaft behandelt und erschöpft, dass nichts zu wünschen übrig bleibt. Wenn man zugestehen darf, dass die von Hrn. \textsc{Bessel} aus einer neuen höchst sorgfältigen Discussion der sämmtlichen guten Beobachtungen abgeleiteten Fundamentalpositionen auf $5''$ in gerader Aufsteigung und Abweichung zuverlässig sind, so wird des Cometen Umlaufszeit zwischen die Grenzen $1691\frac{3}{5}$ und 2673 Jahren eingeschlossen. Eine besondere Aufmerksamkeit verdient Hrn. \textsc{Bessel}'s Verfahren, die Wirkung der Refraction zu berechnen, nur Schade, dass die hiezu gehörigen Formeln durch mehrere Druckfehler ganz entstellt, und daher ohne Zuziehung eines spätern ausführlichen Aufsatzes über diesen Gegenstand in der Mon. Corresp. gar nicht zu gebrauchen sind. 2) Ueber die geographische Länge von Havanna, von Hrn. \textsc{Oltmanns}; nach Beobachtungen von \textsc{Churrucca}, \textsc{Robredo}, \textsc{v. Humboldt} und \textsc{Galiano} im Mittel $5^{\mathrm u}\,38^{\mathrm m}\,51^{\mathrm s}5$ westlich von Paris. 3) Tafeln zur Reduction der Höhen des Polarsterns auf die Meridianhöhe, für die Breite von Berlin. 4) Scheinbare Lichtveränderungen des Algol für die Jahre 1808, 1809, 1810, von Hrn. Prof. \textsc{Wurm}. 5) Drittes und viertes Verzeichniss der verglichenen Lichtstärke der Fixsterne, von Hrn. \textsc{Herschel}. 6) Ueber die wahre geographische Länge des in Peru gemessenen Breitengrades, von Hrn. \textsc{Oltmanns}; aus Beobachtungen von \textsc{Humboldt}'s, \textsc{Ulloa}'s und der Franz. Akademiker im Mittel die Länge von Quito $5^{\mathrm u}\,24^{\mathrm m}\,20^{\mathrm s}$, wofür der letzte Band der Connoissance des tems für 1809 noch $5^{\mathrm u}\,21^{\mathrm m}\,0^{\mathrm s}$ nach \textsc{Bouguer} angab. 7) Tafeln zur Berechnung der jährlichen Veränderung der geraden Aufsteigung und Abweichung der Fixsterne. In diesen Aufsatz haben sich mehrere Unrichtigkeiten eingeschlichen. Nicht die Wirkung der Sonne und des Mondes bringt die jähr\SourceLeaf{511}liche Aenderung $50''15$ in der Länge hervor, die durch die Veränderung der Schiefe der Ekliptik um $0''135$ noch vermehrt wird (wie der Verf. des Aufsatzes sich ausdrückt), sondern die durch die Wirkung der Sonne und des Mondes erzeugte (so genannte Luni-solar-)Präcession wird durch die Verrückung der Ekliptik, die mit der Abnahme der Schiefe zwar einerlei Ursache hat, aber nicht damit verwechselt werden darf, wieder etwas vermindert, und dadurch auf $50''15$, oder nach schärfern Bestimmungen auf $50''10$ heruntergebracht, welches die beobachtete ist. Die Coefficienten des veränderlichen Theils der Präcession in gerader Aufsteigung und die Formel für die Veränderung der Abweichung sollten nicht verschieden, sondern ganz gleich sein; auch die Regel zur Berechnung der Veränderung der jährlichen Variation ist falsch. 8) Beitrag zu den Methoden, eine Reihe Mondsdistanzen für die geographischen Längen in Rechnung zu nehmen, von Hrn. \textsc{Jabbo Oltmanns} (Der Unterschied des scheinbaren und wahren Abstandes wird für drei Zeitmomente berechnet, und daraus für alle Beobachtungen durch Interpolation bestimmt). 9) Astronomische Beobachtungen in Prag 1806, von Hrn. \textsc{David} und \textsc{Bittner}. 10) Methode, durch Hülfe beobachteter Azimuthe, Erhöhungswinkel und relativer Erhöhung irdischer Gegenstände, die geographische Position derselben zu bestimmen, nebst einigen zur Berechnung barometrischer Höhenmessungen dienlichen Hülfstafeln, von Hrn. \textsc{Jabbo Oltmanns}. Viel Genauigkeit kann freilich dies Verfahren nicht geben: indess ist es in ganz unbekannten Gegenden, wenn es an andern Beobachtungen fehlt, nicht zu verwerfen. 11) Ueber das \textsc{Troughton}sche röhrenförmige Pendel, von Hrn. \textsc{Schnitter} in Aachen. In dem Jahrbuche für 1808 war jene \textsc{Troughton}sche Erfindung so beschrieben, dass, sogar nach des Erfinders eigner Angabe für die Ausdehnung des Messings und Stahls, die richtige Compensation nicht herauskam: Hr. \textsc{Schnitter} verdient Dank, hierauf und auf die Nothwendigkeit aufmerksam gemacht zu haben, dass alle abwärts sich ausdehnenden Stangen von Stahl sein müssen, um eine vollkommene Compensation zu erhalten. Inzwischen glaubt Rec., der \textsc{Troughton}'s Originalaufsatz (\textsc{Nicholson}'s philosophical journal. vol. IX.) vor sich hat, bemerken zu müssen, dass der Irrthum blos in der falschen Uebersetzung des Jahrbuchs liegt, und dass \textsc{Troughton}'s Vorschrift gerade so ist, wie sie nach Hrn. \textsc{Schnitter}'s Angabe sein muss. 12) Astronomische Beobachtungen und Bemerkungen von Hrn. \textsc{van Beeck Calkoen}. 13) Beobachtung der Bedeckung von $\alpha$ Scorpion den 2c.\ März 1805, und der Sonnenfinsterniss 16. Juni 1806 auf der Insel Leon bei Cadix, von Hrn. \textsc{Canelas}. 14) Astronomische Beobachtungen zu Wien 1806, von Hrn. \textsc{Triesnecker}. 15) Entdeckung und Beobachtung eines vierten neuen Planeten von Hrn. Dr. \textsc{Olbers}. 16) Beobachtung und Berechnung der Bahn des Cometen von 1806; Beobachtungen der Vesta und Juno, und Sternbedeckungen, von Hrn. F. W. \textsc{Bessel}. 17) Messung der scheinbaren Grösse der Vesta, von Hrn. \textsc{Schröter} (den 26. April 1807 war der scheinbare Durchmesser $0''488$, woraus der wahre Durchmesser 68 Meilen folgt). 18) Beobachtung der Vesta, und Berechnung der Elemente ihrer Bahn und ihres Laufs für 1808, von Hrn. Dr. \textsc{Gauss}. 19) Beobachtungen der Vesta zu Berlin, vom Herausgeber. 20) Astronomische Beobachtungen und Bemerkungen von Hrn. \textsc{Fritsch} in Quedlinburg. Die Beobachtung der Planeten (auch die der Juno im J. 1807, wo sie sehr schwer zu beobachten war) gereichen Hrn. \textsc{Fritsch} zur Ehre, wenn sie schon für den Astronomen wegen Mangels einer gehörigen Zeitbestimmung nicht brauchbar sind. 21) Ueber ein Merkurial-Pendel von Hrn. \textsc{Th. Blacker} in London. 22) Astronomische Beobachtungen von Hrn. \textsc{Lalande}. 23) Beobachtungen über die Climate und Atmosphäre des Saturn von Hrn. Dr. \textsc{Herschel}. Die von Hrn. Dr. H. bemerkten Aenderungen in der Farbe der Polargegenden des Saturn scheinen periodisch zu sein, und Hr. H. sieht sie als einen Beweis vom Dasein einer Atmosphäre an. Astronomische Beobachtungen auf der Berliner Sternwarte 1806; die Kriegsunruhen im Herbste dieses Jahres veranlassten dabei mancherlei Störungen und Unterbrechungen. 25) Lauf der Pallas und Juno im J. 1808 (voraus berechnet). 26) Ueber bemerkte Unterschiede in den scheinbaren Grössen einiger Sterne, von Hrn. Dr. \textsc{Koch} in Danzig. 27) Astronomische Beobachtun\SourceLeaf{512}gen in Spanien (aus einigen Spanischen Schriften gezogen). 28) Einige physisch-astronomische Bemerkungen, von Hrn. \textsc{Huth} in Frankfurt an der Oder. 29) Beitrag zu geographischen Längenbestimmungen, von Hrn. \textsc{Oltmanns}. 30) Vorschlag einer Methode zur Auflösung einer astronomischen Aufgabe, von Hrn. Grafen \textsc{von Platen}, beruht im Wesentlichen auf der Verwandlung von Producten aus Sinus in Summen, deren sich die Astronomen vor Erfindung der Logarithmen wohl bedienten. Jetzt rechnet man in den meisten Fällen bequemer mit den Logarithmen, und gerade in denjenigen, wovon der Verfasser spricht, findet Rec. jenes Verfahren nicht bequem. 31) Vorschlag einer Methode, die Horizontal-Refraction durch die geographische Länge zu bestimmen, von Hrn. \textsc{Oltmanns} (durch Mondsabstände von der niedrig stehenden Sonne). 32) Hrn. \textsc{Fischer}'s Zusatz zu der Abhandlung über die beste Gestalt der Objectiv-Spiegel im Jahrbuch 1806; Hr. F. nimmt hier seine irrige Behauptung, dass sphärische Spiegel den parabolischen vorzuziehen seien, zurück. 33) Verschiedene astronomische Beobachtungen und Nachrichten. Aufmerksamkeit verdient die Anzeige von einem in Remplin von Hrn. \textsc{Hecker} bei $\alpha$ im Füllen den 17. August 1804 beobachteten Stern fünfter Grösse, der nachher nicht wieder gesehen wurde; nur wundern wir uns, dass Hr. \textsc{Hecker} diese Anzeige nicht gleich gemacht hat, wo sie vielleicht von Nutzen hätte sein können.
